\section{Ablation Studies}

\label{sec:ablation}

We conducted ablation experiments on the three primary system components to quantify their individual contributions:

\begin{table}[H]
\centering
\caption{Ablation study results: 12-month retention and field transfer scores.}
\label{tab:ablation}
\begin{tabular}{@{}lcc@{}}
\toprule
\textbf{Configuration} & \textbf{12-Month Retention} & \textbf{Field Transfer} \\
\midrule
Full System & 68.9 & 78.7 \\
$-$ Ecological Grounding & 55.4 ($-$13.5) & 72.4 ($-$6.3) \\
$-$ Adaptive Sequencing & 59.2 ($-$9.7) & 70.1 ($-$8.6) \\
$-$ Multi-Modal Assessment & 64.7 ($-$4.2) & 73.8 ($-$4.9) \\
$-$ All (Static Baseline) & 44.3 ($-$24.6) & 66.6 ($-$12.1) \\
\bottomrule
\end{tabular}
\end{table}

Ecological grounding has the largest individual effect on retention, while adaptive sequencing contributes most to field transfer.
This is consistent with the hypothesis that sensor-triggered reviews strengthen memory consolidation, while personalized learning paths develop deeper procedural competence.

% ============================================================
% 9. Discussion
% ============================================================
